\section{Limits and colimits over preordered sets}
\label{section-posets-limits}

\noindent
A special case of diagrams is given by systems over preordered sets.

\begin{definition}
\label{definition-directed-set}
Let $I$ be a set and let $\leq$ be a binary relation on $I$.
\begin{enumerate}
\item We say $\leq$ is a {\it preorder} if it is
transitive (if $i \leq j$ and $j \leq k$ then $i \leq k$) and
reflexive ($i \leq i$ for all $i \in I$).
\item A {\it preordered set} is a set endowed with a preorder.
\item A {\it directed set} is a preordered set $(I, \leq)$
such that $I$ is not empty and such that $\forall i, j \in I$,
there exists $k \in I$ with $i \leq k, j \leq k$.
\item We say $\leq$ is a {\it partial order} if it is a preorder
which is antisymmetric (if $i \leq j$ and $j \leq i$, then $i = j$).
\item A {\it partially ordered set} is a set endowed with a partial order.
\item A {\it directed partially ordered set} is a directed set
whose ordering is a partial order.
\end{enumerate}
\end{definition}

\noindent
It is customary to drop the $\leq$ from the notation when talking
about preordered sets, that is, one speaks of
the preordered set $I$ rather than of the preordered set $(I, \leq)$.
Given a preordered set $I$ the symbol $\geq$ is defined by
the rule $i \geq j \Leftrightarrow j \leq i$ for all $i, j \in I$.
The phrase ``partially ordered set'' is sometimes abbreviated to ``poset''.

\medskip\noindent
Given a preordered set $I$ we can construct a category: the objects are
the elements of $I$, there is exactly one morphism $i \to i'$
if $i \leq i'$, and otherwise none. Conversely, given a category $\mathcal{C}$
with at most one arrow between any two objects, the set
$\Ob(\mathcal{C})$ is endowed with a preorder defined by the rule
$x \leq y \Leftrightarrow \Mor_\mathcal{C}(x, y) \not = \emptyset$.

\begin{definition}
\label{definition-system-over-poset}
Let $(I, \leq)$ be a preordered set. Let $\mathcal{C}$ be a category.
\begin{enumerate}
\item A {\it system over $I$ in $\mathcal{C}$}, sometimes called a
{\it inductive system over $I$ in $\mathcal{C}$} is given by
objects $M_i$ of $\mathcal{C}$ and for every $i \leq i'$ a
morphism $f_{ii'} : M_i \to M_{i'}$ such that $f_{ii}
= \text{id}$ and such that $f_{ii''} = f_{i'i''} \circ f_{i i'}$
whenever $i \leq i' \leq i''$.
\item An {\it inverse system over $I$ in $\mathcal{C}$},
sometimes called a {\it projective system over $I$ in $\mathcal{C}$}
is given by objects $M_i$ of $\mathcal{C}$ and for every $i' \leq i$ a
morphism $f_{ii'} : M_i \to M_{i'}$ such that $f_{ii}
= \text{id}$ and such that $f_{ii''} = f_{i'i''} \circ f_{i i'}$
whenever $i'' \leq i' \leq i$. (Note reversal of inequalities.)
\end{enumerate}
We will say $(M_i, f_{ii'})$ is a (inverse) system over $I$ to
denote this. The maps $f_{ii'}$ are sometimes
called the {\it transition maps}.
\end{definition}

\noindent
In other words a system over $I$ is just a diagram
$M : \mathcal{I} \to \mathcal{C}$ where $\mathcal{I}$ is the category
we associated to $I$ above: objects are elements of $I$ and
there is a unique arrow $i \to i'$ in $\mathcal{I}$ if and only if $i \leq i'$.
An inverse system is a diagram $M : \mathcal{I}^{opp} \to \mathcal{C}$.
From this point of view we could take (co)limits of any (inverse)
system over $I$. However, it is customary to take
{\it only colimits of systems over $I$} and
{\it only limits of inverse systems over $I$}.
More precisely: Given a system $(M_i, f_{ii'})$
over $I$ the colimit of the system
$(M_i, f_{ii'})$ is defined as
$$
\colim_{i \in I} M_i = \colim_\mathcal{I} M,
$$
i.e., as the colimit of the corresponding diagram.
Given an inverse system $(M_i, f_{ii'})$ over $I$ the limit
of the inverse system $(M_i, f_{ii'})$ is defined as
$$
\lim_{i \in I} M_i = \lim_{\mathcal{I}^{opp}} M,
$$
i.e., as the limit of the corresponding diagram.

\begin{remark}
\label{remark-preorder-versus-partial-order}
Let $I$ be a preordered set. From $I$ we can construct a canonical
partially ordered set $\overline{I}$ and an order preserving map
$\pi : I \to \overline{I}$. Namely, we can define an equivalence
relation $\sim$ on $I$ by the rule
$$
i \sim j \Leftrightarrow (i \leq j\text{ and }j \leq i).
$$
We set $\overline{I} = I/\sim$ and we let $\pi : I \to \overline{I}$
be the quotient map. Finally, $\overline{I}$ comes with a unique
partial ordering such that
$\pi(i) \leq \pi(j) \Leftrightarrow i \leq j$.
Observe that if $I$ is a directed set, then $\overline{I}$
is a directed partially ordered set.
Given an (inverse) system $N$ over $\overline{I}$ we obtain an
(inverse) system $M$ over $I$ by setting $M_i = N_{\pi(i)}$.
This construction defines a functor between the category
of inverse systems over $I$ and $\overline{I}$.
In fact, this is an equivalence.
The reason is that if $i \sim j$, then for any system
$M$ over $I$ the maps $M_i \to M_j$ and $M_j \to M_i$ are
mutually inverse isomorphisms. More precisely, choosing
a section $s : \overline{I} \to I$ of $\pi$ a quasi-inverse
of the functor above sends $M$ to $N$ with
$N_{\overline{i}} = M_{s(\overline{i})}$.
Finally, this correspondence is compatible with colimits of systems:
if $M$ and $N$ are related as above and
if either $\colim_{\overline{I}} N$ or $\colim_I M$ exists
then so does the other and
$\colim_{\overline{I}} N = \colim_I M$.
Similar results hold for inverse systems and limits of inverse systems.
\end{remark}

\noindent
The upshot of Remark \ref{remark-preorder-versus-partial-order}
is that while computing a colimit of a system or a limit of
an inverse system, we may always assume the preorder is a partial order.

\begin{definition}
\label{definition-directed-system}
Let $I$ be a preordered set. We say a system (resp.\ inverse system)
$(M_i, f_{ii'})$ is a
{\it directed system} (resp.\ {\it directed inverse system})
if $I$ is a directed set
(Definition \ref{definition-directed-set}): $I$ is nonempty and
for all $i_1, i_2 \in I$ there exists $i\in I$ such that
$i_1 \leq i$ and $i_2 \leq i$.
\end{definition}

\noindent
In this case the colimit is sometimes (unfortunately)
called the ``direct limit''. We will not use this last
terminology. It turns out that diagrams over a filtered
category are no more general than directed systems in the
following sense.

\begin{lemma}
\label{lemma-directed-category-system}
Let $\mathcal{I}$ be a filtered index category.
There exist a directed set $I$
and a system $(x_i, \varphi_{ii'})$ over $I$ in $\mathcal{I}$
with the following properties:
\begin{enumerate}
\item For every category $\mathcal{C}$ and every diagram
$M : \mathcal{I} \to \mathcal{C}$ with values in $\mathcal{C}$,
denote $(M(x_i), M(\varphi_{ii'}))$
the corresponding system over $I$. If
$\colim_{i \in I} M(x_i)$ exists then so does
$\colim_\mathcal{I} M$ and the transformation
$$
\theta :
\colim_{i \in I} M(x_i)
\longrightarrow
\colim_\mathcal{I} M
$$
of Lemma \ref{lemma-functorial-colimit} is an isomorphism.
\item For every category $\mathcal{C}$ and every diagram
$M : \mathcal{I}^{opp} \to \mathcal{C}$ in $\mathcal{C}$, denote
$(M(x_i), M(\varphi_{ii'}))$ the corresponding inverse system
over $I$. If $\lim_{i \in I} M(x_i)$ exists then so does
$\lim_{\mathcal{I}^{opp}} M$ and the transformation
$$
\theta :
\lim_{\mathcal{I}^{opp}} M
\longrightarrow
\lim_{i \in I} M(x_i)
$$
of Lemma \ref{lemma-functorial-limit} is an isomorphism.
\end{enumerate}
\end{lemma}

\begin{proof}
As explained in the text following
Definition \ref{definition-system-over-poset}, we may view
preordered sets as categories and systems as functors.
Throughout the proof, we will freely shift between these two points of view.
We prove the first statement by constructing a category
$\mathcal{I}_0$, corresponding to a directed set\footnote{In fact,
our construction will produce a directed partially ordered set.}, and a cofinal
functor $M_0 : \mathcal{I}_0 \to \mathcal{I}$. Then, by
Lemma \ref{lemma-cofinal}, the colimit of a diagram
$M : \mathcal{I} \to \mathcal{C}$ coincides with the
colimit of the diagram $M \circ M_0 : \mathcal{I}_0 \to \mathcal{C}$,
from which the statement follows. The second statement is dual to the
first and may be proved by interpreting a limit in $\mathcal{C}$ as
a colimit in $\mathcal{C}^{opp}$. We omit the details.

\medskip\noindent
A category $\mathcal{F}$ is called {\em finitely generated} if
there exists a finite set $F$ of arrows in $\mathcal{F}$, such that
each arrow in $\mathcal{F}$ may be obtained by composing
arrows from $F$. In particular, this implies that $\mathcal{F}$ has
finitely many objects. We start the proof by reducing to the case
when $\mathcal{I}$ has the property that every finitely generated
subcategory of $\mathcal{I}$ may be extended to a finitely
generated subcategory with a unique final object.

\medskip\noindent
Let $\omega$ denote the directed set of finite ordinals, which
we view as a filtered category. It is easy to verify that the
product category $\mathcal{I}\times \omega$ is also filtered,
and the projection
$\Pi : \mathcal{I} \times \omega \to \mathcal{I}$
is cofinal.

\medskip\noindent
Now let $\mathcal{F}$ be any finitely generated
subcategory of $\mathcal{I}\times \omega$.
By using the axioms of a filtered category and a simple induction
argument on a finite set of generators of $\mathcal{F}$,
we may construct a cocone $(\{f_i\}, i_\infty)$ in $\mathcal{I} \times \omega$
for the diagram $\mathcal{F} \to \mathcal{I} \times \omega$. That is, a morphism
$f_i : i \to i_\infty$ for every object $i$ in $\mathcal{F}$
such that for each arrow $f : i \to i'$ in $\mathcal{F}$
we have $f_i = f_{i'} \circ f$. We can also choose $i_\infty$ such
that there are no arrows from $i_\infty$ to an object in $\mathcal{F}$.
This is possible since
we may always post-compose the arrows $f_i$ with an arrow
which is the identity on the $\mathcal{I}$-component and
strictly increasing on the $\omega$-component.
Now let $\mathcal{F}^+$ denote the category consisting of all
objects and arrows in $\mathcal{F}$
together with the object $i_\infty$, the identity
arrow $\text{id}_{i_\infty}$ and the arrows $f_i$.
Since there are no arrows from $i_\infty$ in $\mathcal{F}^+$
to any object of $\mathcal{F}$, the arrow set in $\mathcal{F}^+$
is closed under composition, so $\mathcal{F}^+$ is indeed
a category. By construction, it is a finitely
generated subcategory of $\mathcal{I}$ which has $i_\infty$ as
unique final object. Since, by Lemma \ref{lemma-cofinal},
the colimit of any diagram $M : \mathcal{I} \to \mathcal{C}$
coincides with the colimit of $M\circ\Pi$ , this gives the desired
reduction.

\medskip\noindent
The set of all finitely generated subcategories of $\mathcal{I}$
with a unique final object is naturally ordered by inclusion.
We take $\mathcal{I}_0$ to be the category corresponding
to this set. We also have a functor
$M_0 : \mathcal{I}_0 \to \mathcal{I}$, which takes an
arrow $\mathcal{F} \subset \mathcal{F'}$ in
$\mathcal{I}_0$ to the unique map from the final object of
$\mathcal{F}$ to the final object of $\mathcal{F}'$.
Given any two finitely generated subcategories of
$\mathcal{I}$, the category generated by these two categories is
also finitely generated. By our assumption on $\mathcal{I}$, it is
also contained in a finitely generated subcategory of $\mathcal{I}$
with a unique final object. This shows that $\mathcal{I}_0$ is directed.

\medskip\noindent
Finally, we verify that $M_0$ is cofinal. Since any
object of $\mathcal{I}$ is the final object in the subcategory
consisting of only that object and its identity arrow, the functor
$M_0$ is surjective on objects. In particular, Condition (1) of
Definition \ref{definition-cofinal} is satisfied. Given
an object $i$ of $\mathcal{I}$, objects $\mathcal{F}_1, \mathcal{F}_2$ in
$\mathcal{I}_0$ and maps $\varphi_1 : i \to M_0(\mathcal{F}_1)$
and $\varphi_2 : i \to M_0(\mathcal{F}_2)$ in
$\mathcal{I}$, we can take $\mathcal{F}_{12}$ to be a finitely
generated category with a unique final object containing
$\mathcal{F}_1$, $\mathcal{F}_2$ and the morphisms $\varphi_1, \varphi_2$.
The resulting diagram commutes
$$
\xymatrix{
& M_0(\mathcal{F}_{12}) & \\
M_0(\mathcal{F}_{1}) \ar[ru] & & M_0(\mathcal{F}_{2}) \ar[lu] \\
& i \ar[lu] \ar[ru]
}
$$
since it lives in the category $\mathcal{F}_{12}$ and
$M_0(\mathcal{F}_{12})$ is final in
this category. Hence also Condition (2) is satisfied, which concludes
the proof.
\end{proof}

\begin{remark}
\label{remark-trick-needed}
Note that a finite directed set $(I, \geq)$ always has a greatest object
$i_\infty$. Hence any colimit of a system $(M_i, f_{ii'})$ over such a set
is trivial in the sense that the colimit equals $M_{i_\infty}$. In contrast,
a colimit indexed by a finite filtered category need not
be trivial. For instance, let $\mathcal{I}$ be the category with a single object
$i$ and a single non-trivial morphism $e$ satisfying $e = e \circ e$. The
colimit of a diagram $M : \mathcal{I} \to Sets$ is the image of the
idempotent $M(e)$. This illustrates that something like the trick of passing
to $\mathcal{I}\times \omega$ in the proof of
Lemma \ref{lemma-directed-category-system} is essential.
\end{remark}

\begin{lemma}
\label{lemma-nonempty-limit}
If $S : \mathcal{I} \to \textit{Sets}$ is a cofiltered diagram of sets
and all the $S_i$ are finite nonempty, then $\lim_i S_i$ is nonempty.
In other words, the limit of a directed inverse system of finite nonempty sets
is nonempty.
\end{lemma}

\begin{proof}
The two statements are equivalent by
Lemma \ref{lemma-directed-category-system}.
Let $I$ be a directed set and let $(S_i)_{i \in I}$
be an inverse system of finite nonempty sets over $I$.
Let us say that a {\it subsystem} $T$ is a family $T = (T_i)_{i \in I}$
of nonempty subsets $T_i \subset S_i$ such that $T_{i'}$ is mapped
into $T_i$ by the transition map $S_{i'} \to S_i$ for all $i' \geq i$.
Denote $\mathcal{T}$ the set of subsystems. We order $\mathcal{T}$
by inclusion. Suppose $T_\alpha$, $\alpha \in A$ is a totally ordered family
of elements of $\mathcal{T}$. Say $T_\alpha = (T_{\alpha, i})_{i \in I}$.
Then we can find a lower bound $T = (T_i)_{i \in I}$ by setting
$T_i = \bigcap_{\alpha \in A} T_{\alpha, i}$ which is manifestly a
finite nonempty subset of $S_i$ as all the $T_{\alpha, i}$ are nonempty
and as the $T_\alpha$ form a totally ordered family. Thus we may
apply Zorn's lemma to see that $\mathcal{T}$ has minimal elements.

\medskip\noindent
Let's analyze what a minimal element $T \in \mathcal{T}$ looks like.
First observe that the maps $T_{i'} \to T_i$ are all surjective.
Namely, as $I$ is a directed set and $T_i$ is finite,
the intersection $T'_i = \bigcap_{i' \geq i} \Im(T_{i'} \to T_i)$
is nonempty. Thus $T' = (T'_i)$ is a subsystem contained in $T$ and
by minimality $T' = T$. Finally, we claim that $T_i$ is a singleton
for each $i$. Namely, if $x \in T_i$, then we can define
$T'_{i'} = (T_{i'} \to T_i)^{-1}(\{x\})$ for $i' \geq i$ and
$T'_j = T_j$ if $j \not \geq i$. This is another subsystem as we've seen
above that the transition maps of the subsystem $T$ are surjective.
By minimality we see that $T = T'$ which indeed implies that $T_i$
is a singleton. This holds for every $i \in I$, hence we see that
$T_i = \{x_i\}$ for some $x_i \in S_i$ with $x_{i'} \mapsto x_i$
under the map $S_{i'} \to S_i$ for every $i' \geq i$. In other words,
$(x_i) \in \lim S_i$ and the lemma is proved.
\end{proof}






